\documentclass[12pt,a4paper]{article}
\usepackage{ugmath}
\usepackage{float}
\usepackage{placeins}
\usepackage[labelfont=bf,labelsep=space]{caption}
\newcommand{\paren}[1]{\left( #1 \right)}
\newcommand{\alumno}{Ricardo León Martínez}
\newcommand{\materia}{Fisica I}
\newcommand{\profesor}{Lucero Uscanga Aguilera}
\newcommand{\tarea}{Tarea 5}
\newcommand{\fecha}{13/3/2026}

\begin{document}
    \begin{center}
    {\large\textbf{UNIVERSIDAD DE GUANAJUATO}}\\[0.3cm]
    {\normalsize\textbf{DIVISIÓN DE CIENCIAS NATURALES Y EXACTAS}}\\
    {\normalsize\textbf{CAMPUS GUANAJUATO}}\\[1cm]

    {\Large\textbf{\tarea\ (\materia)}}\\[1cm]
    \end{center}

    \textbf{Nombre:} \alumno \hfill 
    \textbf{Fecha:} \fecha \hfill 
    \textbf{Calificación:} \rule{3cm}{0.4pt} \\[0.3cm]

    \begin{problema}
        Un avión debe volar hacia el norte. La velocidad del avión respecto al aire es
        200km/h y el viento sopla de oeste a este a 90km/h.
        \begin{enumerate}
            \item[(a)] ¿Cuál debe ser el rumbo del avión?
            \item[(b)] ¿Qué velocidad debe llevar el avión respecto al suelo? 
        \end{enumerate}
    \end{problema}

    \begin{problema}
        Dos autos que se desplazan en caminos perpendiculares viajan hacia el norte y el este
        respectivamente. Si sus velocidades con respecto a la tierra son 60km/h y 80km/h, calcular
        la velocidad relativa.
    \end{problema}

    \begin{problema}
        Una granada que se desplaza horizontalmente a una velocidad de 8km/s con respecto a la tierra
        explota en 3 segmentos iguales. Uno de ellos continúa moviéndose horizontalmente a 16km/s,
        otro se desplaza hacia arriba haciendo un ángulo de 45$^{\circ}$ y el tercero se desplaza
        haciendo un ángulo de 45$^{\circ}$ bajo la horizontal. Encontrar la magnitud de las velocidades
        del segundo y tercer fragmentos.
    \end{problema}

    \begin{problema}
        \begin{enumerate}
            \item[(a)] Si el coeficiente de fricción cinética entre neumáticos y pavimento seco es de
            0.80, ¿cuál es la distancia mínima para que se detenga un automóvil que viaja a 28.7m/s bloqueando
            los frenos?
            \item[(b)] En pavimento húmedo, el coeficiente de fricción cinética podría bajar a 0.25.
            ¿Con qué velocidad debemos conducir en pavimento húmedo para poder detenernos en la misma
            distancia que en el inciso (a)?
        \end{enumerate}
    \end{problema}

    \begin{problema}
        La fuerza resultante sobre un objeto de masa $m$ es $F=F_{0}-kt$, donde $F_{0}$ y $k$
        son constantes y $t$ es el tiempo. Encontrar la aceleración. Mediante integración, encontrar
        ecuaciones para la velocidad y la posición.
    \end{problema}

\end{document}